% ============================================================
%  PARTE II: TRADIÇÕES REGIONAIS E CULTURAIS
% ============================================================

\part{TRADIÇÕES REGIONAIS E CULTURAIS}
\partmark{TRADIÇÕES REGIONAIS E CULTURAIS}

% --- Capítulo 5: Helheim - O Reino Gelado dos Mortos na Mitologia Nórdica ---
\chapter[Helheim (Nórdico)]{Helheim - O Reino Gelado dos Mortos na Mitologia Nórdica}
\index{Helheim}
\epigraph{A tradição nórdica apresenta uma das concepções mais singulares de submundo: um reino de gelo e imobilidade, menos suplício físico que aniquilação simbólica, reservado aos que não conquistaram a glória heroica.}

A tradição nórdica apresenta uma das concepções mais singulares de submundo na história das religiões, marcada por um imaginário radicalmente distinto dos modelos mediterrâneos. Enquanto o inferno greco-romano ou cristão enfatiza o fogo purificador e o tormento ativo, o Helheim escandinavo é concebido como um reino de gelo, escuridão e imobilidade, onde a ausência de calor e de vitalidade se converte na própria essência da punição. Trata-se menos de suplício físico e mais de aniquilação simbólica, um destino frio e sombrio reservado àqueles que não conquistaram a glória heroica.

As principais fontes textuais para esse imaginário encontram-se na Edda Prosaica de Snorri Sturluson (c. 1220) e na Edda Poética (compilada entre os séculos XIII-XIV a partir de tradições orais). Entre os poemas, destaca-se a Völuspá, que descreve Helheim como um dos nove mundos suspensos nas raízes de Yggdrasil, a árvore cósmica que sustenta o universo. Essa localização não marginaliza o reino dos mortos como punição divina, mas o integra na ordem cósmica -- Helheim existe porque o equilíbrio do universo exige um espaço para os que morrem de velhice ou doença, em contraposição aos guerreiros conduzidos a Valhalla ou Fólkvangr.

A figura de Hel, filha de Loki e da giganta Angrboda, simboliza essa ambivalência: metade de seu corpo é de carne viva, metade em estado de decomposição. Seu governo não é de maldade absoluta, mas de necessidade cósmica. Helheim é, portanto, menos inferno no sentido cristão do termo e mais um reino paralelo, que preserva o equilíbrio entre vida e morte. A própria Hel reflete a concepção germânica de que o mundo é sustentado pela tensão entre forças opostas, não pela exclusão moral de um lado sobre o outro.


\begin{figure}[ht]
  \centering
  \includegraphics[width=0.85\linewidth]{image21.png}
  \caption{Hel (1889) por Johannes Gehrts – mostra a deusa Hel, metade viva, metade morta, com o cão Garmr aos seus pés, representando o caráter ambíguo e não-moralizante do seu reino.}
  \label{fig:cap5-img1}
\end{figure}

Hel (1889) por Johannes Gehrts -- mostra a deusa Hel, metade viva, metade morta, com o cão Garmr aos seus pés, representando o caráter ambíguo e não-moralizante do seu reino.


\begin{figure}[ht]
  \centering
  \includegraphics[width=0.85\linewidth]{image22.png}
  \caption{Móðguðr guardando a ponte Gjallarbrú (Louis Moe, 1929) – a guardiã impede a passagem até que a alma declare seu nome, simbolizando o limiar entre os vivos e Helheim.}
  \label{fig:cap5-img2}
\end{figure}

Móðguðr guardando a ponte Gjallarbrú (Louis Moe, 1929) -- a guardiã impede a passagem até que a alma declare seu nome, simbolizando o limiar entre os vivos e Helheim.

A geografia mítica de Helheim reforça esse caráter de alteridade. O rio Gjöll, guardado pela donzela Móðguðr, marca a fronteira entre vivos e mortos, atravessada pela ponte Gjallarbrú. O salão de Hel, Éljúdnir, possui teto de ossos de serpentes e piso de entranhas, sugerindo não o tormento ativo, mas um ambiente de decadência permanente e inescapável. O frio, a escuridão e o silêncio são descritos como mais letais do que o fogo -- uma negação da vida em sua essência.

No que tange à punição, a tradição nórdica se afasta de sistemas baseados em julgamento moral universal. A maioria dos mortos em Helheim experimenta uma existência neutra, sem recompensas ou tormentos desmedidos. O sofrimento mais profundo não é o castigo físico, mas a separação da comunidade guerreira, a exclusão da memória heroica e do canto dos escaldos. A vergonha de uma morte não-heroica representa, na ética germânica, uma condenação mais terrível que o suplício corporal.

Uma exceção é Náströnd (“Costa dos Cadáveres”), região afastada do sol, onde perjuros, assassinos e adúlteros sofrem punições explícitas: aprisionados em salões de espinhas de serpentes, são continuamente envenenados por répteis gotejantes. Aqui a tradição nórdica se aproxima de modelos mediterrâneos de inferno punitivo, mas aplica-os apenas a quem violou códigos sociais fundamentais. Para todos os demais, Helheim não é espaço de justiça moral, mas de destinação inevitável.


\begin{figure}[ht]
  \centering
  \includegraphics[width=0.85\linewidth]{image23.jpeg}
  \caption{Visão artística de Helheim como reino frio e sombrio destinado à maioria dos mortos na cosmologia nórdica.}
  \label{fig:cap5-img3}
\end{figure}

O Ragnarök, a batalha final, introduz uma dimensão escatológica própria: mesmo os mortos de Helheim terão papel no fim dos tempos. Conduzidos pelo navio Naglfar, feito de unhas humanas, eles emergirão para lutar contra os deuses. Essa visão contrasta com concepções de condenação estática: na mitologia nórdica, nem mesmo a morte é definitiva -- o cosmos é dinâmico e todos, vivos e mortos, participam de sua destruição e possível renovação.


\begin{figure}[ht]
  \centering
  \includegraphics[width=0.85\linewidth]{image24.jpeg}
  \caption{Naglfar, o navio dos mortos, navegando rumo à batalha final do Ragnarök; ilustração de Louis Moe, 1929.}
  \label{fig:cap5-img4}
\end{figure}

O navio Naglfar é descrito nas Eddas como a embarcação dos condenados, feita de unhas de pessoas mortas; na ilustração de Louis Moe (1929) ele navega rumo ao Ragnarök. O texto do MyNDIR explica que no Prosa Edda o timoneiro é Loki, mas em Völuspá é o gigante Hrymr, e que a cena pertence ao livro Ragnarok: En Billeddigtnin

Com a cristianização da Escandinávia (sécs. X-XII), elementos de Helheim foram reinterpretados. Textos como o Draumkvaede e o Sólarljóð mesclam paisagens gélidas do submundo nórdico com o esquema cristão de julgamento moral, criando um sincretismo fascinante. Ainda assim, a persistência de temas helheimicos em lendas e folclore escandinavos posteriores -- desde aparições fantasmagóricas até crenças em locais amaldiçoados -- demonstra a longevidade cultural dessa concepção singular de pós-morte.


% --- Capítulo 6: Xibalba, Mictlan e Supay - Mundos Subterrâneos das Américas Pré-Colombianas ---
\chapter[Xibalba, Mictlan, Supay]{Xibalba, Mictlan e Supay - Mundos Subterrâneos das Américas Pré-Colombianas}
\index{Xibalba}\index{Mictlan}\index{Supay}


\epigraph{As civilizações pré-colombianas desenvolveram concepções sofisticadas de mundos subterrâneos que rivalizam em complexidade com tradições do Velho Mundo, oferecendo perspectivas únicas sobre punição eterna, julgamento moral e geografia espiritual. Estas tradições, preservadas em códices, monumentos arqueológicos e tradições orais, revelam sistemas cosmológicos integrados onde os submundos constituem componentes essenciais da ordem universal.}

Xibalba, o submundo maia conforme descrito no Popol Vuh (manuscrito k'iche' do séc. XVI baseado em tradições orais antigas), representa uma das mais detalhadas geografias infernais nas Américas. Governado pelos Senhores da Morte, particularmente Hun-Camé ("Um-Morte") e Vucub-Camé ("Sete-Morte"), Xibalba constitui reino de provas e transformações onde almas enfrentam desafios elaborados que determinam seu destino final. A descida dos heróis gêmeos Hunahpú e Ixbalanqué oferece paradigma para navegação pós-morte através de rituais específicos e conhecimento esotérico.


\begin{figure}[ht]
  \centering
  \includegraphics[width=0.85\linewidth]{image25.jpeg}
  \caption{Escultura de Camazotz proveniente de Copán, Honduras; o morcego é associado aos Senhores de Xibalba e aos domínios da noite e da morte.}
  \label{fig:cap6-img1}
\end{figure}

Foto de uma escultura em forma de morcego (Camazotz), proveniente de Copán, Honduras. O morcego (“Killer Bat”) é associado aos Senhores de Xibalba e aos domínios da noite e da morte

A estrutura de Xibalba inclui múltiplas casas de tormento, cada uma representando diferentes aspectos da experiência mortal levados ao extremo. A Casa das Trevas (submersão em escuridão absoluta), Casa das Facas (lâminas cortantes infinitas), Casa do Frio (frieza extrema), Casa dos Jaguares (predadores ferozes), Casa do Fogo (calor consumidor), e Casa dos Morcegos (Camazotz, o deus-morcego decapitador) criam sequência de provas que testam coragem, inteligência e preparação ritual dos mortos.

Mictlan, o submundo asteca, oferece concepção geograficamente estratificada de nove níveis subterrâneos governados por Mictlantecuhtli ("Senhor dos Mortos") e sua consorte Mictlancihuatl ("Senhora dos Mortos"). Diferentemente de sistemas baseados em julgamento moral, Mictlan opera segundo princípios de tipo de morte e preparação ritual. Guerreiros mortos em combate ascendem ao paraíso oriental de Tonatiuhichan, mulheres mortas no parto transformam-se em Cihuateteo (mulheres divinas), enquanto a maioria empreende jornada de quatro anos através dos nove níveis de Mictlan.

A jornada através de Mictlan envolve série de obstáculos geográficos e espirituais que requerem preparação específica. O morto deve atravessar rios caudalosos com ajuda de cão psicopompo (xoloitzcuintli), superar montanhas que se chocam, resistir a ventos cortantes capazes de despir carne dos ossos, e enfrentar serpentes e lagartos

gigantescos. Esta jornada culmina na audiência com Mictlantecuhtli, onde o morto apresenta oferendas específicas (jade, incenso, tecidos) para garantir descanso final.


\begin{figure}[ht]
  \centering
  \includegraphics[width=0.85\linewidth]{image26.jpeg}
  \caption{Mictlantecuhtli e Quetzalcóatl na página 56 do Códice Borgia, associados à morte e à renovação da vida.}
  \label{fig:cap6-img2}
\end{figure}

Página 56 do Códice Borgia: o deus da morte Mictlantecuhtli aparece à esquerda e Quetzalcóatl à direita; juntos representam a primeira morte e a segunda vidacommons.wikimedia.org. A obra é uma reprodução fidedigna do manuscrito pré-colombiano, em domínio público.

A concepção inca de Supay e o submundo andino reflete adaptação a geografia montanhosa única. Supay, frequentemente traduzido como "diabo" pelos missionários espanhóis, representa na verdade divindade ctônica complexa que governa tanto riquezas minerais quanto mortos. O submundo inca localiza-se nas profundezas terrestres, acessível através de cavernas e fontes naturais consideradas portais entre mundos. Esta geografia sagrada integra paisagem física andina com cosmologia espiritual, criando mapa sagrado where specific locations serve as interfaces entre worlds.

Os processos de colonização europeia e evangelização cristã nas Américas criaram sincretismos complexos entre tradições pré-colombianas e conceitos cristãos de inferno. A interpretação de Xibalba como "inferno cristão" pelos missionários facilitou preservação de elementos indígenas sob terminologia cristã. Documentos coloniais como os Anales de Cuauhtitlan e o Códice Florentino preservam descrições de submundos pré-colombianos filtradas através de conceitos cristãos, criando sínteses únicas.

A persistência de elementos pré-colombianos em tradições populares contemporâneas demonstra vitalidade cultural destas concepções. O Día de los Muertos mexicano preserva elementos de comunicação com mortos em Mictlan, enquanto tradições andinas mantêm oferendas a Supay para proteção em atividades mineiras. Estas continuidades culturais ilustram adaptabilidade de cosmologias indígenas face a pressões evangelizadoras europeias.


\section{Comparação entre Xibalba, Mictlan e Ukhu Pacha}

\begin{description}[leftmargin=2.8cm,style=nextline]
  \item[Xibalba] Governado por Hun-Camé e Vucub-Camé, os Senhores da Morte.
  Suas casas de tormento submetem os viajantes a provas de trevas, lâminas,
  frio, jaguares, fogo e morcegos.
  \item[Mictlan] Governado por Mictlantecuhtli e Mictlancihuatl. A travessia
  compreende nove níveis e exige que o morto supere sucessivos obstáculos.
  \item[Ukhu Pacha] Mundo inferior andino governado por Supay, divindade
  ctônica associada aos mortos e às riquezas minerais.
\end{description}

As casas de tormento de Xibalba incluem:

\begin{itemize}
  \item Trevas
  \item Facas
  \item Frio
  \item Jaguares
  \item Fogo
  \item Morcegos (Camazotz)
\end{itemize}
Os nove níveis da travessia de Mictlan incluem:

\begin{itemize}
  \item Rio com cão psicopompo
  \item Montanhas colidindo
  \item Ventos cortantes
  \item Serpentes
  \item Lagartos gigantes
  \item Obstáculos adicionais
\end{itemize}
O Ukhu Pacha é um submundo ligado a cavernas, minas e fontes naturais.

Em Xibalba, o acesso depende de provas rituais e conhecimento esotérico;
seu horizonte narrativo é a transcendência ou dissolução, representada
pelos heróis gêmeos Hunahpú e Ixbalanqué. Em Mictlan, o destino depende do
tipo de morte e da preparação ritual: após quatro anos de travessia e das
oferendas a Mictlantecuhtli, alcança-se o repouso. No mundo andino, o acesso
é universal e não moralista, ligado à integração com os ancestrais e com as
forças da terra.

A análise comparativa revela paralelos estruturais notáveis entre submundos americanos e tradições do Velho Mundo, sugerindo arquétipos recorrentes na organização de conceitos escatológicos. Todos apresentam geografia estratificada, provas rituais, divindades governantes duais (masculina/feminina) e ênfase em preparação adequada para a navegação pós-morte. Essas similaridades, desenvolvidas independentemente, sugerem constantes antropológicas na conceituação humana do destino pós-morte.


\begin{figure}[ht]
  \centering
  \includegraphics[width=0.85\linewidth]{image27.jpeg}
  \caption{Máscara boliviana de China Supay em gesso policromado, início do século XX, pertencente ao Musée du quai Branly.}
  \label{fig:cap6-img3}
\end{figure}

Máscara de China Supay, em gesso policromado, originária da Bolívia (início do século XX). A peça faz parte das coleções do Musée du quai Branly e representa o ser ctônico que preside o submundo andino e as riquezas minerais.


% --- Capítulo 7: Tradições Célticas e Germânicas ---
\chapter[Tradições Célticas]{Tradições Célticas e Germânicas}
\index{Annwn|see{Outro Mundo céltico}}\index{Outro Mundo céltico}\index{Sídhe}


\epigraph{As tradições célticas e germânicas pré-cristãs desenvolveram concepções distintivas de mundos subterrâneos e punição pós-morte que exerceram influência profunda na formação do cristianismo medieval europeu. Estas tradições, preservadas em literatura vernácula, folclore e evidências arqueológicas, revelam sistemas cosmológicos sofisticados que resistiram à cristianização através de processos de sincretismo e reinterpretação cultural.}

O conceito céltico de Annwn, o Outro Mundo subterrâneo descrito particularmente na literatura galesa medieval, representa um reino paralelo acessível através de lagos, colinas sagradas (sidhe) e outros portais naturais. O Mabinogi, coletânea de contos galeses dos sécs. XI-XIII baseados em tradições orais mais antigas, apresenta Annwn como reino governado por reis como Arawn e Gwyn ap Nudd, onde diferentes leis temporais e morais operam. Significativamente, Annwn não constitui local de punição, mas reino alternativo onde heróis podem obter objetos mágicos e conhecimento esotérico.

A literatura irlandesa preserva concepções similares através do conceito de Sídhe, tanto como locais físicos (montículos funerários neolíticos) quanto como dimensão espiritual habitada pelos Tuatha Dé Danann após sua retirada do mundo humano. O Táin Bó Cúailnge e outros ciclos épicos irlandeses descrevem interações frequentes entre mundos humano e sidhe, onde heróis como Cú Chulainn empreendem jornadas ao Outro Mundo para completar geas (obrigações místicas) ou obter armas sobrenaturais.

A tradição germânica continental, parcialmente preservada em textos como o Heliand saxão (séc. IX) e evidências arqueológicas de práticas funerárias, revela concepções de submundo relacionadas mas distintas das tradições nórdicas. O conceito de Hel germânico primitivo aparece associado não apenas à divindade feminina dos mortos, mas a estados de ocultação e transformação espiritual. Evidências arqueológicas de "boat burials" e oferendas mortuárias sugerem crenças em jornadas pós-morte através de reinos aquáticos subterrâneos.

Os processos de cristianização das Ilhas Britânicas (sécs. V-VIII) criaram sínteses únicas entre tradições célticas e conceitos cristãos de inferno. A conversão de reis como Æthelberht de Kent (c. 601) e Edwin da Nortúmbria (627) facilitou desenvolvimento de literatura cristã vernácula que preserva elementos païos significativos. O poema anglo- saxão Christ and Satan (séc. VIII-IX) apresenta descrições de inferno que incorporam elementos de paisagens nórdicas e germânicas, incluindo ênfase em frieza e desolação.

A influência druídica na formação de conceitos cristãos irlandeses representa fenômeno de particular importância. Textos como o Acallamh na Senórach (séc. XII) preservam tradições sobre druidas convertidos que mantêm conhecimento esotérico sobre navegação entre mundos. A figura de Merlim na literatura arturiana representa síntese de sabedoria druídica com elementos cristãos, criando arquétipo de sábio que compreende tanto Annwn quanto inferno cristão.

A literatura visionária irlandesa medieval, particularmente textos como Fís Adamnáin (Visão de Adamnán, séc. X-XI) e Aisling Óenguso (Sonho de Óengus), desenvolve geografias de Outro Mundo que integram elementos célticos tradicionais com cosmologia cristã. Estas visões apresentam paisagens que combinam sidhe tradicional com descrições de purgatório e inferno, criando sínteses geográficas únicas que influenciarão desenvolvimentos continentais posteriores.

O folclore céltico posterior preserva elementos pré-cristãos através de contos sobre "fairy hills", "ghost roads", e outros fenômenos liminais que conectam mundo dos vivos com reinos dos mortos. A persistência de tradições como Samhain (Halloween) demonstra continuidade cultural de conceitos sobre comunicação temporária entre mundos durante períodos específicos do calendário astronômico.

A análise linguística de terminologias célticas e germânicas para conceitos escatológicos revela stratificação cultural complexa. Termos irlandeses como tech duinn ("casa escura") e aos sí ("folk dos sidhe") preservam elementos de cosmologia pré-cristã, enquanto empréstimos latinos como peccad (pecado) e diabul (diabo) demonstram influência cristã. Esta coexistência terminológica reflete sincretismo conceitual profundo que caracteriza cristianismo céltico.


% --- Capítulo 8: Tradições Eslavas e Bizantinas ---
\chapter[Tradições Eslavas]{Tradições Eslavas e Bizantinas}
\index{Nav}


\epigraph{As tradições eslavas e bizantinas representam desenvolvimentos únicos na evolução dos conceitos infernais, combinando elementos de paganismo eslavo primitivo com sofisticadas elaborações teológicas orientais. Esta síntese, facilitada pela conversão do mundo eslavo ao cristianismo bizantino entre os sécs. IX-XI, criou tradições escatológicas distintivas que contrastam significativamente com desenvolvimentos occidentais contemporâneos.}

O conceito eslavo primitivo de Nav representa um dos mais antigos estratos de pensamento escatológico eslavo. Nav constitui dimensão paralela onde residem mortos, mas também forças espirituais ativas que influenciam mundo dos vivos. Diferentemente de concepções de submundo como reino separado, Nav permeia realidade ordinária como camada oculta acessível através de rituais específicos, sonhos, e estados alterados de consciência. Esta permeabilidade fundamental distingue cosmologia eslava de modelos mediterrâneos de separação radical entre mundos.

A documentação de crenças eslavas pré-cristãs apresenta desafios significativos devido à transmissão principalmente oral e aos esforços sistemáticos de cristianização. Fontes como a Crônica Primária (séc. XII) e o Slovo o polku Igoreve (séc. XII) preservam elementos de mitologia eslava filtrados através de perspectivas cristãs. Divindades como Veles/Volos, associadas ao submundo e aos mortos, aparecem demonizadas em textos cristãos, mas mantêm características que revelam concepções pagãs originais.

A conversão da Rússia de Kiev sob Vladimir I (988 d.C.) marca transformação fundamental nas concepções escatológicas eslavas. A adoção do cristianismo bizantino introduz teologia sofisticada sobre inferno, purgatório e salvação, mas estes conceitos interagem com substratos pagãos persistentes. A Palavra sobre Lei e Graça de Hilarião de Kiev (c. 1050) representa uma das primeiras sínteses sistemáticas de teologia bizantina com sensibilidades eslavas, enfatizando temas de renovação cósmica e redenção universal.

O desenvolvimento da theosis (deificação) na tradição bizantina oferece perspectiva escatológica distintiva que influencia concepções eslavas sobre destino pós-morte. Padres como Máximo o Confessor (580-662) e João Damasceno (676-749) elaboram doutrinas onde salvação constitui participação na natureza divina ao invés de simples escape da punição. Esta ênfase em transformação positiva ao invés de punição negativa cria atmosfera teológica diferente da que predomina no Ocidente latino.

A literatura apocalíptica eslava eclesiástica, incluindo textos como o Apocalipse de Abraão eslavo e a Escatologia de Metodio, desenvolve visões escatológicas que integram elementos bizantinos com sensibilidades eslavas específicas. Estes textos, frequentemente pseudepígrafos baseados em originais gregos, apresentam descrições detalhadas de inferno e julgamento final que incorporam elementos de geografia eslava e preocupações culturais específicas.

O monasticismo ortodoxo eslavo, particularmente através de figuras como Teodósio de Kiev (1008-1074) e Sérgio de Radonez (1314-1392), desenvolve espiritualidade ascética que enfatiza preparação para morte através de práticas hesiquiastas. A oração de Jesus e outras técnicas contemplativas criam experiência antecipatória de estados pós-morte, oferecendo alternativa experiencial a especulações teóricas sobre inferno.

A iconografia bizantina eslava desenvolve representações visuais distintivas de escatologia através de afrescos e ícones que decoram igrejas e mosteiros. A Última Julgamento de Novgorod (séc. XV) e frescos similares em Pskov e Moscou apresentam geografia infernal que combina elementos bizantinos tradicionais com motivos eslavos locais, incluindo paisagens de inverno e arquitectura vernácula.

O folclore eslavo posterior preserva elementos pré-cristãos através de contos sobre rusalkas (espíritos aquáticos), domovoi (espíritos domésticos), e outras entidades liminais que habitam fronteiras entre mundo dos vivos e dos mortos. A persistência destas tradições em contextos ostensivamente cristãos demonstra sincretismo profundo que caracteriza cristianismo eslavo popular.

As divergências entre tradições orientais e ocidentais tornam-se evidentes durante debates sobre Filioque e outras questões que eventualmente culminarão no Grande Cisma (1054 d.C.). Diferenças sobre pneumatologia refletem abordagens distintas à autoridade teológica que afetam interpretações escatológicas. Tradições orientais mantêm maior abertura à teologia apofática e mistérios escatológicos, enquanto desenvolvimentos ocidentais enfatizam sistematização jurídica e definições precisas sobre estados pós-morte.

O legado dos Grandes Concílios para desenvolvimentos escatológicos posteriores é fundamental. As definições doutrinárias estabelecidas durante estes séculos criam framework teológico que determina interpretações subsequentes de conceitos infernais tanto no Oriente quanto no Ocidente. A cristalização de ortodoxia trinitária e cristológica elimina alternativas especulativas que poderiam ter levado a concepções radicalmente diferentes sobre punição eterna, estabelecendo parâmetros doutrinários que persistem até o presente.
